% Proposed chapter for the Picat User's Guide, Version 4.0.
% Insert after predfunc.tex. The anonymous forms describe the proposed extension;
% parser support alone does not provide the compiler/runtime implementation.
\chapter{Higher-order Calls and Anonymous Definitions}
\label{chapter:higher-order}

A higher-order predicate or function takes a callable value as an argument.
Picat supports callable values that name predicates or functions, optionally
with arguments already supplied. Anonymous functions and anonymous predicates
provide callable values without requiring named definitions.

\section{Higher-order Call Conventions}
\index{higher-order call}
\label{section:higher-order-conventions}

The predicate \texttt{call($S$, $Arg_1$, $\ldots$, $Arg_n$)} invokes a
predicate, while the function \texttt{apply($S$, $Arg_1$, $\ldots$, $Arg_n$)}
returns the value of a function call. For a named callable, $S$ is an atom or
a structure. If $S$ is an atom, it supplies the name of the predicate or
function. If $S$ is a structure, its name supplies the callable name and its
arguments precede the additional arguments. For example:
\begin{verbatim}
F = $'*'(2),
Y = apply(F,5)              % Y = 10
\end{verbatim}
The resulting function call is \texttt{'*'(2,5)}. The dollar symbol constructs
a structure instead of evaluating a function call. A bare atom can also name
a callable:
\begin{verbatim}
F = '+',
Y = apply(F,3,4)            % Y = 7
\end{verbatim}

When $S$ is a variable at compile time, both call conventions can be written
as \texttt{$S$($Arg_1$, $\ldots$, $Arg_n$)}. The context determines whether
this notation denotes a predicate call or a function call:
\begin{verbatim}
P = $member(X),
P([1,2,3]),                % predicate context: call(P,[1,2,3])
F = $'*'(2),
Y = F(5)                   % function context: apply(F,5)
\end{verbatim}
An occurrence used as a goal invokes a predicate; an occurrence used as a
value-producing expression invokes a function. The result variable of a
function call is not an additional explicit argument.

Within a term constructor, a higher-order call whose name is a variable is
also treated as a function call:
\begin{verbatim}
F = $'*'(2),
T = $p(F(5))               % T = p(10)
\end{verbatim}
Here, the value of \texttt{F(5)} becomes an argument of the constructed term.
The higher-order call is not preserved as literal syntax in that term.

\subsection{Arity and Supplying Arguments}
The supplied arguments must select a predicate or function of the appropriate
arity. In particular, \texttt{apply} does not return a partially evaluated
function when too few arguments are supplied; an incorrect arity raises an
exception. A structure such as \texttt{\$'*'(2)} explicitly stores a supplied
argument. Constructing that structure is different from invoking a function
with too few arguments.

An anonymous definition has the explicit arity given by its parameter list.
All of those arguments must be supplied when it is invoked. To return a
function that takes further arguments, write a nested anonymous function
explicitly, as illustrated below.

\section{Anonymous Functions}
\index{anonymous function}\index{lambda expression}
\label{section:anonymous-functions}

An anonymous function, also called a lambda expression, has the form
\begin{verbatim}
([A1,A2,...,An] -> Res, B1,B2,...,Bm)
\end{verbatim}
The outer parentheses are required. The bracketed list contains the parameter
patterns, \texttt{Res} is the return expression, and the goals following the
first comma form the body. The body is executed before the return expression
is evaluated. When no body is needed, the function can be written as
\begin{verbatim}
([A1,A2,...,An] -> Res)
\end{verbatim}
Omitting the body is equivalent to using the body \texttt{true}.

Evaluating a lambda expression constructs a callable value, called a
\emph{closure}; it does not execute the body or evaluate the return expression.
No dollar symbol is needed to construct this value. For example:
\begin{verbatim}
F = ([X,Y] -> X+Y),
Z = F(3,4)                 % Z = 7
\end{verbatim}
The invocation \texttt{F(3,4)} can also be written as
\texttt{apply(F,3,4)}. A function with a body can compute its result using
intermediate variables:
\begin{verbatim}
F = ([X,Y] -> Z, T = X*X, Z = T+Y),
R = F(3,4)                 % R = 13
\end{verbatim}
The body has the usual goal syntax, including conjunctions, disjunctions,
conditionals, and loops. For example:
\begin{verbatim}
Sum = ([Xs] -> Total,
    Total = 0,
    foreach (X in Xs)
        Total := Total+X
    end),
S = Sum([1,2,3])           % S = 6
\end{verbatim}
The function follows the execution conventions of an ordinary named function.
It is not a generator of multiple return values.

\subsection{Parameter Patterns}
The parameters are terms, just like the arguments of a named rule head; they
are not restricted to variables. Their matching follows the ordinary head
pattern-matching conventions. For example:
\begin{verbatim}
First = ([[X|_]] -> X),
H = First([a,b,c])         % H = a
\end{verbatim}
The outer bracketed list contains one parameter, whose pattern is a nonempty
list. Matching does not instantiate the caller's argument merely to make a
pattern applicable. The same restrictions on head patterns apply here;
parameter patterns are not evaluated as function calls.

A parameter list must be a proper list of patterns. Its length determines the
arity. An empty list denotes a function with no explicit arguments:
\begin{verbatim}
F = ([] -> 42),
Y = F()                    % Y = 42
\end{verbatim}

\section{Anonymous Predicates}
\index{anonymous predicate}
\label{section:anonymous-predicates}

An anonymous predicate has the form
\begin{verbatim}
([A1,A2,...,An] => B1,B2,...,Bm)
\end{verbatim}
Its parameter list has the same interpretation as that of an anonymous
function. The body is a goal, and there is no return expression. For example:
\begin{verbatim}
Print = ([X] => println(X)),
Print(hello)
\end{verbatim}
The invocation \texttt{Print(hello)} can also be written as
\texttt{call(Print,hello)}. An anonymous predicate with no explicit arguments
uses \texttt{[]} as its parameter list and is invoked with \texttt{P()} or
\texttt{call(P)}.

An anonymous predicate consists of one rule, written with \texttt{=>}; there
is no anonymous \texttt{?=>} form. There is no alternative rule to try on
backtracking. This does not require every goal in the body to be globally
deterministic: a goal such as \texttt{member} can still supply multiple
solutions within the body. For example:
\begin{verbatim}
Choose = ([X] => member(X,[a,b,c])),
L = findall(X,Choose(X))   % L = [a,b,c]
\end{verbatim}
Use an anonymous predicate when success, failure, or bindings are the desired
result; use an anonymous function when an expression must return a value.

\section{Context Variables and Nested Definitions}
\index{closure}\index{context variable}\index{name scope}

An anonymous definition forms a \emph{name scope}, just like a
\texttt{foreach} loop. A variable that has occurred before the anonymous
definition in the enclosing scope is \emph{global} to the anonymous
definition and refers to the existing variable. A variable whose first
occurrence is inside the anonymous definition is a fresh \emph{local}
variable of that definition. This rule applies to variables in parameter
patterns, the return expression, and the body. Here, ``global'' describes
visibility across the scope boundary; it does not mean a module-global
variable. A prior occurrence counts even if the variable is still unbound.

Local variables are fresh for each invocation and are not made visible to
subsequent goals outside the definition. For example:
\begin{verbatim}
Offset = 10,
F = ([X] -> R, T = X+Offset, R = T*T),
A = F(1),                 % A = 121
B = F(2)                  % B = 144
\end{verbatim}
The variable \texttt{Offset} is global to the anonymous definition, while
\texttt{X}, \texttt{R}, and \texttt{T} first occur inside it and are local.
The two invocations use fresh local variables. Nested anonymous definitions
follow the same scope rule: a variable introduced in an outer definition can
be global to an inner definition.

An anonymous definition can refer to variables from the enclosing context.
A closure carries the values of these context variables at the time the
closure is created. Subsequent changes to the context variables do not affect
the values captured by that closure. This capture convention applies to both
anonymous functions and anonymous predicates. For example:
\begin{verbatim}
Offset = 10,
AddOffset = ([X] -> X+Offset),
Offset := 20,
Y = AddOffset(5),         % Y = 15
NewAddOffset = ([X] -> X+Offset),
Z = NewAddOffset(5)       % Z = 25
\end{verbatim}
The parameter \texttt{X} is supplied at invocation, while \texttt{Offset}
comes from the context in which the closure is constructed. The captured
variables are not additional arguments that the caller must supply.
The first closure retains the captured value \texttt{10}, even after
\texttt{Offset} is assigned \texttt{20}. The second closure is created after
the assignment and therefore captures \texttt{20}. Invoking either closure
does not recapture the current value of \texttt{Offset}.

An anonymous function can return another anonymous function:
\begin{verbatim}
MakeAdder = ([A] -> ([B] -> A+B)),
Add10 = MakeAdder(10),
Y = Add10(5)               % Y = 15
\end{verbatim}
The returned closure retains the context needed to use \texttt{A} after the
outer invocation finishes. This example uses two explicit invocations rather
than an implicit partial application.

\section{Passing Anonymous Definitions as Arguments}
A higher-order routine can invoke an anonymous definition using the same
conventions as a named callable. For example:
\begin{verbatim}
transform(F,[]) = [].
transform(F,[X|Xs]) = [F(X)|transform(F,Xs)].

visit(P,[]) => true.
visit(P,[X|Xs]) =>
    P(X),
    visit(P,Xs).
\end{verbatim}
These definitions can be used as follows:
\begin{verbatim}
Ys = transform(([X] -> X*X),[1,2,3]),
visit(([X] => println(X)),Ys)
\end{verbatim}
The transformation returns \texttt{[1,4,9]}, and the visit prints those values.
Neither routine needs to know whether its argument was constructed from a
named callable or an anonymous definition.

\section{Modules and Call Resolution}
\index{higher-order call!module resolution}
For an ordinary named higher-order call whose target is not known at compile
time, the runtime system searches the implicitly imported built-in modules
\texttt{basic}, \texttt{io}, \texttt{math}, and \texttt{sys}, followed by the
other loaded modules in the order returned by \texttt{loaded\_modules()}, and
then the global module. If the compiler knows the callable name, it binds the
call using the ordinary compile-time resolution conventions.

A variable call can be module-qualified, as in \texttt{m.F(X,Y)}. For this
form, lookup starts in \texttt{m}, followed by the implicitly imported
built-in modules and the global module. The module qualifier is an atom,
not a variable. The restriction on constructing literal dot-notation terms
still applies; qualification is call syntax rather than a module-qualified
structure that can be constructed as data. Ordinary named higher-order terms
cannot reference private symbols, which are compiled away at compile time.

A closure is associated with its defining code and captured context. Passing
it to another module must preserve that association: calls written inside an
anonymous body use the defining module's resolution context. The receiving
module can invoke the closure without importing the defining module.

\section{Implementation of Anonymous Definitions}
\index{anonymous definition!implementation}
The compiler introduces a unique dummy name for each anonymous definition
and compiles the definition just like a named function or predicate. The
dummy name belongs to the global module. The body, however, is compiled in
the module in which the anonymous definition occurs. Thus, calls in the
body are resolved using the defining module's definitions and imports,
including its private definitions, just as in an ordinary named definition
in that module. The return expression of an anonymous function is compiled
in the same defining context.

The closure carries the generated name together with the values of the
captured context variables. Conceptually, an anonymous function such as
\begin{verbatim}
Offset = 10,
F = ([X] -> X+Offset)
\end{verbatim}
is replaced by a closure of the form
\begin{verbatim}
F = $lam999([Offset])
\end{verbatim}
and the compiler introduces the function
\begin{verbatim}
lam999([Offset],X) = X+Offset => true.
\end{verbatim}
Here, \texttt{lam999} stands for a unique compiler-generated name, and the
first argument carries the captured environment. Invoking \texttt{F(5)}
supplies the explicit argument after that environment argument. An anonymous
predicate is translated in the same way, using a generated predicate rather
than a function. For example, the anonymous predicate
\begin{verbatim}
([X] => println(X+Offset))
\end{verbatim}
has a corresponding generated definition of the form
\begin{verbatim}
lam1000([Offset],X) => println(X+Offset).
\end{verbatim}
Its closure is conceptually \texttt{\$lam1000([Offset])}.

Placing the generated name in the global module allows a receiving module
to invoke the closure without importing the defining module. Compiling the
body in the defining module preserves the meaning of calls within the body.
The generated name is unique, so separately compiled anonymous definitions
do not share a dummy name. The names and environment lists shown here
illustrate the translation; programs should not depend on their particular
spelling or representation.

\section{Compiled Bodies and Callable Terms}
Anonymous bodies are compiled code. They can contain assignments, loops, and
conditionals, even though these constructs cannot be preserved as literal
callable terms for general \texttt{call} interpretation. An anonymous
predicate is therefore useful for passing such a computation to another
routine.

The printed representation and generated name of a closure are implementation
details. Programs should construct, pass, and invoke the closure rather than
inspect or reconstruct its generated representation. No particular runtime
lookup strategy or performance guarantee is implied by the anonymous syntax.
